% !TeX root = ../geometry_main.tex
\documentclass[../geometry_main.tex]{subfiles}
\begin{document}
%==================================================================
%======================================================================
\chapter{Curvature \& torsion}\label{sec:curvature_torsion}
%========================================================================
%======================================================================

%-----------------------------------------------------------------------
%======================================================================
\section{Torsion \& the equivalence principle}
\label{sec:torsion_and_equivalence_principle}
%========================================================================
%-------------------------------------------------------------------------

\begin{mytheorem}[Torsion measures the nonclosure of infinitesimal parallelograms]
Consider a manifold \(M\) with affine connection \(\nabla\).
Let \(X,Y\) be two vector fields on \(M\).
Construct a \emph{parallelogram} with vertex \(p \in M\) and parallel axes $X(p), Y(p)$ by first flowing in direction \(X(p)\) and then in the direction parallel to the other axis $Y(p)$, or viceversa.
For infinitesimal distances, the first curve is approximated by $\gamma_{X}^\epsilon\simeq p+\epsilon X(p)$, but the direction parallel to the other axis is $Y(p+\epsilon X)+\epsilon \nabla_X Y(p)$, and not simply $Y(p+\epsilon X)$, since by definition of parallel transport along $\gamma_X$ 
{\small\begin{align}
    \nabla_X Y(p)=\lim_{\epsilon \to 0}\frac{Y(p+\epsilon X)-\tau^\epsilon_{\gamma_X}(Y(p))}{\epsilon}\,.
\end{align}}
The two orders give different possible endpoints, namely 
{\small\begin{align}
    &p\mapsto p + \epsilon X(p)\mapsto p + \epsilon X(p) + \epsilon \big[Y(p+\epsilon X) + \epsilon \nabla_X Y(p)\big]= p+ \epsilon \big[X(p) + Y(p)\big] + \epsilon^2\Big[X^\mu\partial_\mu Y(p)+ \nabla_X Y(p)\Big],
    \\[3pt]
    &p\mapsto p + \epsilon Y(p)\mapsto p + \epsilon Y(p) + \epsilon \big[X(p+\epsilon Y) + \epsilon \nabla_Y X(p)\big]= p+ \epsilon \big[X(p) + Y(p)\big] + \epsilon^2\Big[Y^\mu\partial_\mu X(p)+ \nabla_Y X(p)\Big].
\end{align}}
The non-closure of the would-be (infinitesimal) parallelogram is measured by the torsion tensor 
\begin{equation}
    \Delta p = \epsilon^2 T(X,Y) + \mathcal{O}(\epsilon^3), \qquad T(X,Y) := \nabla_X Y - \nabla_Y X - [X,Y],
\end{equation}
Beware that \emph{finite} parallelograms in general do not close in curved spaces even in absence of torsion\footnote{Just think of a parallelogram on the surface of an apple}, but infinitesimal ones do close in familiar contexts.
The presence of torsion hence measures this somewhat puzzling aspect.
In fact, one can avoid this counterintuitive geometric picture and recast torsion as a `fifth force' violating the equivalence principle, so that the interpreation becomes that infinitesimal parallelograms do close, but the presence of an extra force field makes particles deviate from the geodesic.
\end{mytheorem}

\begin{figure}[htbp]
    \centering
    \includegraphics[width=0.62\textwidth]{images/torsion_parallelogram.pdf}
    \caption{Torsion and nonclosure of infinitesimal parallelograms.}
    \label{fig:torsion-parallelogram}
\end{figure}


\begin{mytheorem}[Why torsion is a tensor?]
The fact that torsion is a tensor, even though the Christoffel symbols are not, follows directly from its geometric definition as measuring the nonclosure of parallelograms.
Alternatively, we could view it from the `physicitsts definition of tensors' and its antisymmetric nature.
Under a change of frame the Christoffel symbols transform as
\begin{equation}
    \Gamma^{\rho}_{\mu \nu} = \frac{\partial x^{\rho}}{\partial x^{\prime \sigma}} \frac{\partial x^{\prime \alpha}}{\partial x^{\mu}} \frac{\partial x^{\prime \beta}}{\partial x^{\nu}} \Gamma^{\prime \sigma}_{\alpha \beta} + \frac{\partial x^{\rho}}{\partial x^{\prime \sigma}} \frac{\partial^2 x^{\prime \sigma}}{\partial x^{\mu} \partial x^{\nu}}.
\end{equation}
The nontensorial part of the transformation is symmetric in the lower indices \(\mu, \nu\).
Therefore by splitting
\begin{equation}
    \Gamma^{\rho}_{\mu \nu} = \Gamma^{\rho}_{(\mu \nu)} + \Gamma^{\rho}_{[\mu \nu]},
\end{equation}
we see that indeed $T_{\mu,\nu}^{\quad \rho} = 2 \Gamma^{\rho}_{[\mu \nu]}$ transforms as a tensor.
\end{mytheorem}


\begin{mytheorem}[Weak, standard and strong equivalence principle]
There are several inequivalent statements usually referred to as `the equivalence principle', listed below in increasing order of strength.
Notably, it is possible to \textbf{obey the weak and standard equivalence principle while violating the strong one with fifth forces} (cf. arXiv:2204.08484v4, 	arXiv:2309.11496v2).
{\scriptsize
\begin{center}
\setlength{\tabcolsep}{3pt}
\renewcommand{\arraystretch}{1.15}
\begin{tabular}{c|p{0.25\textwidth}|p{0.25\textwidth}|p{0.25\textwidth}}
\textbf{Equivalence Principle} & \textbf{Statement} & \textbf{Applies to} & \textbf{Excludes/tests} \\
\hline
Weak
&
Freely falling test bodies follow the same trajectories independently of composition and internal structure: \(m_{\rm inertial}=m_{\rm gravitational}\).
&
Non-self-gravitating matter,  treated as test particles.
&
Composition-dependent fifth forces; violations of universality of free fall.
\\
\hline
Standard
&
Locally in a freely falling lab, all non-gravitational laws reduce to special relativity, independently of velocity and spacetime location.
&
All local non-gravitational physics: mechanical experiments, electromagnetism, nuclear physics, atomic clocks, etc.
&
Violations of local Lorentz invariance; non-metric couplings of matter.
\\
\hline
Strong
&
The above local-free-fall statement also applies to gravitational physics, and self-gravitating bodies fall universally.
&
Matter and gravity itself, including planets, stars, black holes and gravitational waves.
&
Scalar-tensor effects; varying \(G\); Nordtvedt effect; dependence on gravitational binding energy.
\end{tabular}
\end{center}}
%
Non-self-gravitating bodies are object whose internal structure is \emph{not} determined by gravity, but rather by other forces of the standard model. That is object that are viriliazed by forces other than gravity, like electromagnetic or nuclear forces.
For example, an atom, a pen and a car are non-self-gravitating bodies, while a star, a balck hole, a galaxy and a galaxy cluster are self-gravitating bodies.
\end{mytheorem}



\begin{mytheorem}[Torsion \& the equivalence principle]
As far as the metric is concerned, the \emph{standard} equivalence principle is \emph{completely equivalent} to the vanishing of torsion.
Indeed, the \emph{weak} equivalence principle states that at any point there exists a \emph{local freely falling frame}, i.e. a frame where the symmetric part of the Christoffel symbols vanish at that point.
Physically, this demands the geodesic equation reduce to a straight line
\begin{align}
    \nabla_{\dot{x}}\dot{x} = \frac{d^2 x^{\sigma}}{d \lambda^2} + \Gamma^{\sigma}_{\mu \nu} \frac{d x^{\mu}}{d \lambda} \frac{d x^{\nu}}{d \lambda} = 0
    \quad \xrightarrow[\text{local free fall at $p$}]{} \quad
    \frac{d^2 x^{\sigma}}{d \lambda^2}_{\mid p} = 0.
\end{align} 
Imposing this holds for any choice of initial velocity \(\frac{d x^{\mu}}{d \lambda}_{\mid p}\) implies the symmetric part $\Gamma_{(\mu\nu)}^\sigma=0$ at \(p\).

Note that this seems to leave room for torsion--i.e. for nonvanishing antisymmetric part $\Gamma_{[\mu\nu]}^\sigma$.
However, the \emph{standard} equivalence principle actually demands \emph{anynon-gravitational experiment} reduce to the special relativistic one in such local freely falling frames, including mechanical ones.
In particular, it requires that in absence of forces the angular momentum $L$ of spinning bodies stays constant, i.e. it is parallely transported via the connection.
That is 
\begin{align}
    0=\nabla_{\dot{x}}L = \frac{d L^\sigma}{d \lambda} + \Gamma^{\sigma}_{\mu \nu} \frac{d x^\mu}{d \lambda} L^\nu 
    \quad \xrightarrow[\text{local free fall at $p$}]{} \quad
    \frac{d L^\sigma}{d \lambda}_{\mid p} = 0.
\end{align}
Since both vectors $L=L^\sigma \partial_\sigma$ and $\tfrac{dx^\sigma}{d \lambda}\partial_\sigma$ are arbitrary at \(p\), all of $\Gamma_{\mu\nu}^\sigma$ must vanish at \(p\).

In fact, as we mentioned, the vanishing ot torsion not only is needed, but is completely \emph{equivalent} to the existence of such local freely falling frames as we show here below.
\end{mytheorem}


\begin{theorem}
Let \(M\) be a manifold with an affine connection \(\nabla\), not necessarily metric-compatible.
The torsion of \(\nabla\) vanishes at a point \(p \in M\) if and only if there exists a local coordinate system \(\{x^{\mu}\}\) around \(p\) such that the connection coefficients vanish \(\Gamma^{\rho}_{\mu \nu} (p) = 0\) at $p$.
\end{theorem}
\begin{proof}
Since torsion is a tensor, if it vanishes in one frame then it vanishes in all frames and is thus identically zero.
Conversely, Riemann normal coordinates--i.e. coordinates obtained via the exponential map at \(p\)--always exist and by construction they are such that the \emph{symmetric} part $\Gamma^{\rho}_{(\mu \nu)}(p)=0$ vanishes at this point.
Indeed, in such cohordinates stratight lines $\lambda\mapsto \exp(\lambda v)$ are mapped to geodesics through $p$ in $M$ 
\begin{equation}
    \graycancel{\frac{d^2}{d \lambda^2} (\lambda v^\sigma)}+ \Gamma^{\sigma}_{\mu \nu} v^\mu v^\nu = 0.
\end{equation}
Since $v\in T_pM$ is arbitrary, this implies $\Gamma^{\sigma}_{(\mu \nu)}(p)=0$.
\end{proof}


\begin{mytheorem}[Quantum gravity \& allowing for nonzero torsion]
If we want to quantize gravity, it is wise and natural to allow for some nonzero torsion.
Indeed, in quantum mechanics all quantities are subject to quantum fluctuations.
Therefore, if the metric is quantum fluctuating, it might well be the case that it explores also configurations with nonzero torsion.

In fact, we have several examples where quantum fuctuations violate classical constraints.
An example is quantum tunnelling, where  particles can cross classically forbidden regions.
Another one is the nonvanishing of correlators between spacelike separated operators in quantum field theory, violating classical causality.
In this second case, the `rules' of classical general relativity are restored by the fact that these a-causal correlators are exponentially suppressed at scales larger than the Compton wavelength of the particles involved.
Furthermore, the analogous contributions from particles and antiparticles cancel each other in the actually observable correlator.
\end{mytheorem}



\begin{mytheorem}[Torsion may not be forbidden, just Wilson-suppressed]
The reason why we do not see torsion at accessible energies might be different, and deeper in a sense, than just advocating the equivalence principle. 
Instead of viewing torsion only as forbidden by the equivalence principle, we can view it as an additional field i.e. a possible fifth force.
Namely, given a metric-compatible connection with torsion, we split it into Levi-Civita connection and torsion tensor as $\Gamma^\rho_{\mu\nu} = \bar{\Gamma}^{\rho}_{\mu\nu} + T^{\rho}_{\mu\nu}$.
The effect of torsion is then moved to the right-hand side of Einstein equations
\begin{equation}
    \bar{G}_{\mu\nu}=8\pi G \left(T_{\mu\nu}^{\rm matter} + T_{\mu\nu}^{\rm torsion}\right),
\end{equation}
and interpreted as extra matter-like sources, or as fifth-force interactions

In the Wilsonian framework however, unless protected by a symmetry, such a field will be integrated out at low energies.
Indeed, naive power-counting suggests that having 3-indices torsion will generate non-renormizable interactions, much more suppressed by some heavy scale \(M_T\) at low energies than the 2-indices metric interactions.
\end{mytheorem}



\begin{mytheorem}[The Levi-Civita torsion-fixed metric connection.]
Given a pseudoriemannian manifold \((M,g)\) and a $(1,2)$-tensor $T$ antisymmetric in the lower indices, there exists a unique connection \(\nabla\) that is metric-compatible and has torsion \(T\), i.e. 
\begin{equation}
    \nabla g = 0,  \qquad T_\nabla(X,Y):=\nabla_XY-\nabla_YX-[X,Y]\equiv T(X,Y).
\end{equation}
The (anti)-symmetric parts $\Gamma_{\mu\nu}^\rho=\Gamma^\rho_{(\mu\nu)}+\Gamma^\rho_{[\mu\nu]}$ of the Christoffel symbols are uniquely determined by 
\begin{equation}\label{eq:levi_civita_christoffel_symbols}
    \Gamma^\rho_{(\mu\nu)}= \frac{1}{2} g^{\rho \sigma} \left( \partial_{\mu} g_{\nu \sigma} + \partial_{\nu} g_{\mu \sigma} - \partial_{\sigma} g_{\mu \nu} \right), \qquad 
    \Gamma^\rho_{[\mu\nu]}= T_{\mu \nu}^{\rho}.
\end{equation}
%
The \textbf{contraction of the (symmetrized) Christoffel symbols} of the Levi-Civita connection satisfy the identity
\begin{equation}\label{eq:metric_christoffel_contraction_identity}
    \Gamma^{\mu}_{(\mu \nu)} = \frac{1}{\sqrt{|g|}} \partial_{\nu} \sqrt{|g|}.
\end{equation}   
\end{mytheorem}
\begin{proof}
From the Levi-Civita formula \eqref{eq:levi_civita_christoffel_symbols} we have
\begin{equation}
    \Gamma^{\mu}_{(\mu \nu)} = \frac{1}{2} g^{\mu \sigma} \left( \partial_{\mu} g_{\nu \sigma} + \partial_{\nu} g_{\mu \sigma} - \partial_{\sigma} g_{\mu \nu} \right)
    = \frac{1}{2} g^{\mu \sigma} \partial_{\nu} g_{\mu \sigma} + \underbrace{\frac{1}{2} g^{\mu \sigma} \left( \partial_{\mu} g_{\nu \sigma} - \partial_{\sigma} g_{\mu \nu} \right)}_{=0},
\end{equation}
where two terms cancel due to the symmetry of \(g^{\mu \sigma}\) and dummy indices.
Differentiating the determinant of the metric, 
\begin{align}
\begin{aligned}
    \partial_{\nu} \sqrt{|g|} &= \frac{1}{2 \sqrt{|g|}} \text{sign}(g) \det(g_{\mu \nu}) \, \text{Tr}\big(g^{\alpha \beta} \partial_{\nu} g_{\beta \gamma}\big)
    = \frac{1}{2} \sqrt{|g|}\,\, g^{\alpha \beta} \partial_{\nu} g_{\beta \alpha}
    = \sqrt{|g|} \,\,\Gamma^{\mu}_{(\mu \nu)}.
\end{aligned}
\end{align}
\end{proof}


\begin{mytheorem}[Interpretation of the Christoffel symbols]
In view of what discssed above, we can interpret the Christoffel symbols and their derivatives as follows.
\begin{itemize}
    \item The antisymmetric part of the Christoffel symbols-- the torsion--measures the presence of fifth forces fields that violate the equivalence principle i.e. beyond standard gravity.
    Indeed, we showed that nonzero torsion is equivalent to the impossibility of finding a local freely falling frame.
    This antisymmetric part can \emph{not} be absorbed into (the symmetric part of) the energy-momentum tensor, since this must indeed match the \emph{symmetric} Einstein tensor.
    The presence of torsion thus genuinely corresponds to `extra forces' beyond standard gravity.
    %
    \item The symmetric part of the Christoffel symbols corresponds to `standard gravity' and `standard pseudoforces', like Coriolis force, that can be eliminated by going to a freely falling frame.
    Indeed nonzero Christoffel symbols also appears in flat spacetime when using curved coordinates, like spherical coordinates.

    Notably, the choice of what to attribute to gravity and what to inertial forces is \emph{conventional} and depends on the choice of connection.
    Indeed, the difference of two connections $\nabla,\tilde{\nabla}$ is indeed a $(1,2)$-tensor since Christoffel symbols transform as
    \begin{align}\label{eq:difference_connection_is_tensor}
        \Gamma^{\rho}_{\mu \nu}- \tilde{\Gamma}^{\rho}_{\mu \nu} = \frac{\partial x^{\rho}}{\partial x^{\prime \sigma}} \frac{\partial x^{\prime \alpha}}{\partial x^{\mu}} \frac{\partial x^{\prime \beta}}{\partial x^{\nu}} \left(\Gamma^{\prime \sigma}_{\alpha \beta}-\tilde{\Gamma}^{\prime \sigma}_{\alpha \beta}\right) + \graycancel{ \frac{\partial x^{\rho}}{\partial x^{\prime \sigma}} \frac{\partial^2 x^{\prime \sigma}}{\partial x^{\mu} \partial x^{\nu}}}.
    \end{align}
    If both of them are symmetric in the lower indices, their difference is also symmetric and can thus be interpreted as a field contributing to the energy-momentum tensor at the rhs of Einstein's equations.
    %
    \item Finally, derivatives of Christoffel symbols measure curvature, that is tidal forces that cannot be eliminated even when switching to a freely falling frame.
    Indeed, a freely falling frame is `exactly freely falling' only at a single point, but as soon as we depart from it infinitesimal tidal forces start to act.
    Curvature--i.e. essentially the Riemann tensor $R\sim \partial \Gamma$--therefore measures how much nearby freely falling frames deviate from each other.
\end{itemize}
\end{mytheorem}


%-----------------------------------------------------------------------
%======================================================================
\section{Curvature}\todotag{To be written}
%========================================================================
%---------------------------------------------------------------------------


%======================================================================
\subsection{What is curvature \& how to quantize it?}
%========================================================================


\begin{mytheorem}[How to measure curvature? How to quantize gravity?]
%
What is curvature? How do we even measure it? Intuitively, curvature is how much a space deviates from being flat.
Here are some ways to quantify this.
In order to quantize gravity--i.e. curvature--it will be useful to have several notions of curvature, since we do not know yet which one will be the most appropriate for a quantum theory of gravity and what object will be quantized.
To appreciate this point, note that we are used to measuring space-time with rulers and clocks, but already at scales smaller than an atom these notions lose meaning, without even invoking the Planck length. 


\begin{enumerate}
\item \textbf{Violation of angle sum laws.}  For a triangle with angles \(\alpha,\beta,\gamma\) one has
\[
\alpha + \beta + \gamma \neq 180^\circ.
\]
Such a deficit angle can be used to encode the shape. This notion of curvature thus ultimately boils down to \emph{measuring angles and possibly quantize them}. (This idea is ideed used in some approaches to quantum gravity.)
%
\item \textbf{Violation of Pythagoras' law.}  For a triangle with side lengths \(a,b,c\) one has
\[
a^2 + b^2 \neq c^2.
\]
In this case we \emph{encode the shape through metric distances}--i.e. via a Riemannian metric \(g\) on \(M\)--and we might \emph{quantize the metric} $g$ itself (as in canonical quantum gravity) or some functionals of it (as in loop quantum gravity).
%
\item \textbf{Tensorial definiton i.e. tidal forces} (Riemann curvature tensor and Gaussian curvature).
 This is an umbrella term for the most common notion of curvature, where this is encoded in suitable tensors.
The most common is the Riemann tensor $R$, but we might equivalently define it via other \emph{equivalent} tensorial means like the \textbf{sectional curvature}.
We would then quantize a suitable expression of these?
%
\item \textbf{Nontrivial parallel transport around loops.}  A vector transported around a closed loop generally does not return to itself.  One can therefore encode the shape through an affine connection \(\Gamma\) on \(M\), which we might then quantize somehow.
%
\item \textbf{Green functions of quantum fields i.e. nontrivial entanglement of vacuum fluctuations.} Namely, n-point correlators of quantum fields do know about the geometry of the underlying manifold.
For example, the 2-point function of a free scalar field theory on a curved manifold \((M,g)\) satisfies
\begin{equation}
    (\Box_x + m^2) G(x,y) = - \frac{1}{\sqrt{|g|}} \delta^{(n)} (x-y),
\end{equation}
where \(\Box\) is the D'Alambert operator associated to the metric \(g\).
From the knowledge of the Green function \(G(x,y)\) for all pairs of points \((x,y) \in M \times M\), one might try to reconstruct the metric \(g\) and thus the curvature of the manifold. 
In doing so, quantization conditions on $g$ might emerge naturally from the quantum nature of the field theory itself.
%
\item \textbf{Spectral geometry i.e. non-triviality of the sound of the manifold.}  The spectrum of differential operators defined on a manifold \((M,g)\) encodes the geometry of the manifold itself.
The idea is that the spectrum of differential operators defined on a manifold \((M,g)\) encodes the geometry of the manifold itself.
For example, the eigenvalues of the Laplace-D'Alambert operator \(\Box\) depend on the metric \(g\) and thus on the curvature of \((M,g)\).
One might then try to quantize gravity by quantizing the spectrum of such operators. (cf. relevant section) 
\end{enumerate}
%
Crucially, the above \textbf{local descriptions carry redundant information}  (excluding Green functions and spectral geometry).
Indeed two pseudo-Riemannian manifolds \((M,g)\) and \((M,g')\) should be considered equivalent— they describe the same spacetime— if there exists an isometric diffeomorphism
\[
e:(M,g)\longrightarrow (M,g'), \quad \text{s.t.}\quad g = e^*(g').
\]
In other words \(e\) is merely a change of chart, the two manifolds are just different coordinate descriptions of the \emph{same} manifold and hence have the same \emph{shape}.
This redundancy makes it hard to identify the true degrees of freedom that should be quantized.
This is the gauge-fixing problem already well-known from other gauge theories like QED.
%
\end{mytheorem}



\begin{mytheorem}[Pseudo-Riemannian structures \& fixing the gauge.\notetag{Polish}]
%
If \((M,g)\) and \((M,g')\) are related by an isometry then they are merely different coordinate descriptions of the same manifold and hence have the same `shape'
%
A \textbf{pseudo-Riemannian structure} \(\mathcal{E}\) is an equivalence class of pseudo-Riemannian manifolds which can be mapped into each other via metric-preserving diffeomorphisms i.e. coordinate changes.
Hence space(time) should be modelled as a (pseudo)-Riemannian structure \(\mathcal{E}\), i.e. an equivalence class of pairs \((M,g)\).

One would like to \textbf{fix the gauge} i.e. identify exactly one representative \((M,g)\) for each equivalence class \(\mathcal{E}\).
These \textbf{classes} are \textbf{hard} to handle \textbf{because the nonlinearity of geometry} makes it hard to establish the (in)existence of an isometry between two generic curved metrics \(g\) and \(g'\).
The flat metric (class) is easy to identify, regardless of complicated coordinates, since it is equivalent to the vanishing of the Riemann tensor, but there is no known general algorithm to establish the isometry of two arbitrary metrics.
Instead one typically relies on (incomplete) sets of invariants, like curvature scalars, holonomies, spectra of differential operators, K-theory classes, etc.

This issue is particularly evident in the \textbf{path integral formulation of gravity}.
Indeed we would like to integrate over all possible \emph{inequivalent} structures \(\mathcal{E}\)
\begin{equation}
    \int_{\mathrm{geom. structs.}} e^{i S[\mathcal{E}]} \;D\mathcal{E},
\end{equation}
for a suitable isometry-invariant action $S[M,g]\equiv S[\mathcal{E}]$, like the Einstein-Hilbert action.
%
What we initially have is however an integral over all metrics or connections or other objects uniquely identifying a riemannian structure, i.e.
\begin{equation}
    \int_{\mathrm{metrics}} \!\!Dg\, e^{i S[g]} 
    \quad \text{or} \quad
    \int_{\mathrm{connections}}\!\!D\Gamma\,  e^{i S[\Gamma]}\quad \text{etc.}
\end{equation}
%
We must then introduce a gauge-fixing functional \(\delta(\mathcal{E})\) so that from each equivalence class of metrics, or other objects uniquely identifying a pseudo-Riemannian structure, only one representative contributes to the path integral.
%
A viable option is to perform the integration directly over \textbf{non-redundant invariants} that uniquely identify a pseudo-Riemannian structure, like spectra of suitable differential operators on tensor valued forms or K-theory classes, so that no gauge-fixing is needed.

How many  actual physical degrees of freedomhas a pseudo-Riemannian structure in \(n\) dimensions?
A symmetric covariant 2-tensor \(g_{\mu \nu}\) has \(\frac{n(n+1)}{2}\) independent components at each point, but diffeomorphisms eliminate \(n\) of them, which are spurious artifacts of coordinates.
The \textbf{number of physical degrees of freedom of a pseudo-Riemannian structure} is thus \(\frac{n(n-1)}{2}\).
In \(n=4\) dimensions we have 6 physical DoFS at each point.
Notably, in $n=2$ we only have 1 DoF, confirming that curvature in 2D is completely specified by a single scalar, the Gaussian curvature i.e. the sectional curvature of the only 2D plane in the tangent space.
In particular any 2d metric is \emph{locally}\footnote{There might be global topological obstructions.} conformally flat i.e. can be written as \(g_{\mu\nu}=e^{\phi(x)}\eta_{\mu\nu}\) for some scalar function \(\phi\).
Also note for $n=1$ there are no DoFs: there is no curvature in 1 dimension.
%
\end{mytheorem}


\begin{mytheorem}[Weyl \& Ricci curvature decomposition\todotag{Finish}]
%
In $n>3$ the Ricci tensor does not encode all the information of the Riemann tensor i.e. all the DoFs of the curvature, the remaining being encoded in the Weyl tensor.
In general, the Riemann curvature tensor \(R_{\alpha \beta, \mu \nu}\) can be decomposed into the \textbf{Weyl tensor} \(C_{\alpha \beta, \mu \nu}\) and terms depending on the Ricci tensor \(R_{\mu \nu}\) and the Ricci scalar \(R\) as\todotag{Check}
\begin{equation}
    R_{\alpha \beta, \mu \nu} = C_{\alpha \beta, \mu \nu} + \frac{1}{n-2} \left( g_{\alpha [\mu} R_{\nu] \beta} - g_{\beta [\mu} R_{\nu] \alpha} \right) - \frac{R}{(n-1)(n-2)} g_{\alpha [\mu} g_{\nu] \beta}.
\end{equation}
The Weyl tensor is traceless and encodes the conformal part of the curvature, while the Ricci part encodes the volume-changing part of the curvature.
Notably the Weyl tensor can be nonzero in vacuum spacetimes and accounts for gravitational waves, whereas the Ricci tensor vanishes in the absence of matter.

According to the power counting above, in 2D the metric has only 1 DoF, completely encoded in the Ricci scalar, and the Weyl tensor vanishes identically.
In 3D the metric has 3 DoFs, the (symmetric) Ricci tensor also has $3=\frac{3(3+1)}{2}-3$ independent components, and the Weyl tensor can be expressed entirely in terms of the Ricci tensor and scalar.
%
\end{mytheorem}




%=======================================================
\subsection{The Riemann curvature tensor}
%=======================================================

The Riemann curvature tensor is defined as the operator measuring the noncommutativity of covariant derivatives.

\begin{mytheorem}[Ricci identity: the Riemann tensor measures noncommutativity of parallel transport in two directions]\todotag{finish}
If we neglect torsion the path obtained by first following the flow of a vector field \(X\) for infinitesimal distance \(\epsilon\), then the flow of a vector field \(Y\) for infinitesimal distance \(\epsilon\), then the flow of \(-X\) for infinitesimal distance \(\epsilon\), and finally the flow of \(-Y\) for infinitesimal distance \(\epsilon\), closes to leading order in \(\epsilon\).
In presence of torsion, the parallelogram does not close, but the Riemann tensor retains its meaning as measuring the noncommutativity of covariant derivatives.
\end{mytheorem}


\begin{figure}[htbp]
    \centering
    \includegraphics[width=0.62\textwidth]{images/curvature_parallelogram.pdf}
    \caption{Curvature and noncommutativity of parallel transport around infinitesimal parallelograms.}
    \label{fig:curvature-parallelogram}
\end{figure}

\begin{mytheorem}[The Riemann tensor is just a rotation matrix in disguise\todotag{Finish}]
For the sake of the interpretation, assume torsion vanishes, everything holds even if not.
For fixed $X,Y\in \mathfrak{X}(M)$, the Riemann curvature tensor \(R(X,Y)\) is a linear map
\begin{equation}
    R(X,Y): T_p M \to T_p M \quad \mid Z\mapsto R(X,Y)Z,
\end{equation}
that returns the vector obtained by parallel transporting \(Z\) infinitesimally first in the $X$ direction and then in the $Y$ direction, and subtracting the vector obtained by doing the same but first in the $Y$ direction and then in the $X$ direction.
This is nothing but an infinitesimal rotation of the vector \(Z\) around the infinitesimal parallelogram spanned by \(X\) and \(Y\).
The same interpretation holds in presence of torsion, but the parallelogram does not close anymore and we have to correct for that.
In this sense, at fixed $X$ and $Y$ the Riemann tensor is just a rotation matrix in disguise.
%
\end{mytheorem}




\begin{mytheorem}[Symmetries of the Riemann tensor\todotag{Finish}]
The Riemann curvature tensor \(R\) has the following symmetries.
\begin{itemize}
    \item \textbf{Antisymmetry in last two indices} (no need metric, even if torsion is present):
    \begin{equation}
        R(X,Y)(\cdot) = - R(Y,X)(\cdot) \quad \Longrightarrow \quad
        R^\alpha_{\,\beta,\, \mu\nu} = - R^\alpha_{\,\beta,\, \nu\mu}.
    \end{equation}
    \item \textbf{Antisymmetry in first two lowered indices} (if metric compatible, \emph{even} if torsion is present):
    \begin{equation}
        g(R(X,Y)Z,W) = - g(R(X,Y)W,Z) \quad \Longrightarrow \quad
        R_{\alpha \beta,\, \mu\nu} = - R_{\beta \alpha,\, \mu\nu}.
    \end{equation}
    \item \textbf{Symmetry under exchange of pairs of indices} (if metric compatible) \todotag{not sure if torsion is present}
    \begin{equation}
        R_{\alpha \beta,\, \mu\nu} = R_{\mu \nu,\, \alpha \beta}.
    \end{equation}
    \item \textbf{First Bianchi identity}
    \begin{align}
        R(X,Y)Z& + R(Y,Z)X + R(Z,X)Y = \nabla_X T(Y,Z) + \nabla_Y T(Z,X) + \nabla_Z T(X,Y) \\
        & + T(T(X,Y),Z) + T(T(Y,Z),X) + T(T(Z,X),Y).
    \end{align}
    In components, this reads
    \begin{equation}       
        R^\alpha_{\,\beta,\, \mu\nu} + R^\alpha_{\,\mu,\, \nu\beta} + R^\alpha_{\,\nu,\, \beta\mu} = \dots
    \end{equation}
    \item \textbf{Second Bianchi identity}
    \begin{equation}
        \nabla_X R(Y,Z) + \nabla_Y R(Z,X) + \nabla_Z R(X,Y) = \dots
    \end{equation}
    In components, this reads
    \begin{equation}
        \nabla_\lambda R^\alpha_{\,\beta,\, \mu\nu} + \nabla_\mu R^\alpha_{\,\beta,\, \nu\lambda} + \nabla_\nu R^\alpha_{\,\beta,\, \lambda\mu} = \dots
    \end{equation}
\end{itemize}
    
\end{mytheorem}
\begin{proof}
The first identity just follows from the definition of the Riemann tensor 
\begin{align}
    R(X,Y) &= \nabla_X \nabla_Y - \nabla_Y \nabla_X - \nabla_{[X,Y]} \\
    & = - \left( \nabla_Y \nabla_X - \nabla_X \nabla_Y - \nabla_{[Y,X]} \right) = - R(Y,X).
\end{align}
The second follow from metric compatibility
\begin{align}
    \nabla g=0 &\Rightarrow X \big( g(Z,W) \big) = \graycancel{(\nabla_X g)(Z,W)}+ g(\nabla_X Z, W) + g(Z, \nabla_X W).
\end{align}
Then 
\begin{align}
    g(R(X,Y)Z,W) &= g(\nabla_X \nabla_Y Z, W) - g(\nabla_Y \nabla_X Z, W) - g(\nabla_{[X,Y]} Z, W)
    \\
    & = X \big( g(\nabla_Y Z, W) \big) - Y \big( g(\nabla_X Z, W) \big) - [X,Y]( g(Z, W)) \\
    & - g(\nabla_Y Z, \nabla_X W) + g(\nabla_X Z, \nabla_Y W) + g(Z, \nabla_{[X,Y]} W) \\
    & = \dots
\end{align}

\end{proof}



\begin{mytheorem}[Sectional curvature]\todotag{Finish.}
Sectional curvature is the simplest and most familiar notion of curvature one can define on a Riemannian manifold \((M,g)\).
Given a point \(p \in M\) and a 2-dimensional subspace \(\sigma \subseteq T_p M\) of the tangent space at \(p\), the sectional curvature \(K(\sigma)\) is defined in various equivalent ways.
\begin{itemize}
    \item \textbf{Gauss curvature.}  It is defined as the Gaussian curvature of the surface obtained by exponentiating \(\sigma\) via the exponential map at \(p\).
    Namely, given the surface $S = \{\exp_p (v) \,|\, v \in \sigma \}$, we set
    \begin{equation}
        K(\sigma) := K_{Gauss} (S).
    \end{equation}
    \item \textbf{Riemann tensor.}  We define
    \begin{equation}
        K(\sigma) := \frac{R(u,v,v,u)}{||u||^2 ||v||^2 - (g(u,v))^2},
    \end{equation}
    where \(\{u,v\}\) is any basis of \(\sigma\).
\end{itemize}
\end{mytheorem}


\begin{mytheorem}[Einstein manifolds\todotag{write}]

\end{mytheorem}


\begin{mytheorem}[Kulkarni--Nomizu tensor in constant curvature spaces\todotag{finish}]
In a (pseudo-)Riemannian manifold of constant sectional curvature \(K\), the Riemann curvature tensor can be expressed as
\[
    R_{abcd} = K (g_{ac} g_{bd} - g_{ad} g_{bc}).
\]
This is proved using the properties of the Kulkarni--Nomizu product and the fact that sectional curvature uniquely determines the Riemann tensor.
%
\end{mytheorem}
\end{document}
